% 背景・課題
\section{背景・課題}

\subsection{組み込み開発における一般的な課題}

組み込みシステムの開発において、以下のような問題が頻繁に発生する：

\begin{enumerate}
    \item \textbf{delay()の乱用} \\
    ブロッキング関数である\texttt{delay()}を使用すると、その間は他の処理が完全に停止する。
    複数のセンサーやアクチュエータを同時に制御する場合、\texttt{delay()}は致命的なボトルネックとなる。

    \item \textbf{モジュール間の密結合} \\
    各ハードウェアモジュールの制御コードが\texttt{loop()}内に直接記述されると、コードの肥大化・複雑化が進み、
    モジュールの追加・削除・変更が困難になる。

    \item \textbf{出力の重複・不整合} \\
    統一的なライフサイクル管理がない場合、同一フレーム内で複数回のセンサー読み取りや出力が
    発生するなど、処理の重複や不整合が起こりやすい。

    \item \textbf{グローバル変数の氾濫} \\
    モジュール間のデータ共有にグローバル変数を多用すると、依存関係が不明瞭になり、
    バグの原因となる。
\end{enumerate}

\subsection{ESP32-S3 CAMの特徴}

本プロジェクトで使用するESP32-S3 N16R8 CAM開発ボードは、以下の特徴を持つ：

\begin{table}[h]
\centering
\begin{tabularx}{0.8\textwidth}{l X}
\toprule
\textbf{項目} & \textbf{仕様} \\
\midrule
MCU & ESP32-S3（デュアルコア Xtensa LX7） \\
Flash & 16MB（N16） \\
PSRAM & 8MB（R8、OPI接続） \\
通信 & WiFi + Bluetooth 5.0 \\
カメラ & OV2640 / OV5640 対応 \\
フレームワーク & Arduino \\
\bottomrule
\end{tabularx}
\caption{ESP32-S3 N16R8 CAM 主要スペック}
\end{table}

多機能なハードウェアを効率的に管理するために、統一的なモジュール管理アーキテクチャが求められる。
